\section{Objectives and Methodology}

\begin{frame}{Research Objectives}
    \begin{block}{Our goal}
        Turn the modulo trick from a useful idea into a precise method we can test, prove, and apply.
    \end{block}

    \vspace{0.15cm}
    \begin{columns}[T,onlytextwidth]
        \column{0.48\textwidth}
        \begin{exampleblock}{1. Define}
            \small Solvable and non-solvable by the modulo trick.
        \end{exampleblock}

        \begin{exampleblock}{2. Build}
            \small A simple framework for the method.
        \end{exampleblock}

        \column{0.48\textwidth}
        \begin{exampleblock}{3. Prove}
            \small When the method fails.
        \end{exampleblock}

        \begin{exampleblock}{4. Validate}
            \small Brenner and Foster observations.
        \end{exampleblock}
    \end{columns}
\end{frame}

\begin{frame}{Research Methodology}
    \begin{figure}[h!]
        \centering
        \scalebox{0.5}{\begin{tikzpicture}[
            node distance=0.8cm,
            every node/.style={font=\small},
            box/.style={
                rectangle,
                draw,
                rounded corners,
                align=center,
                minimum width=4.5cm,
                minimum height=0.8cm,
                fill=blue!10
            },
            arrow/.style={->, thick}
        ]
        
        % Nodes
        \node[box] (1) {Proposition 1: Upper bound $\implies$ finite solutions};
        \node[box, below=of 1] (2) {Define solvable/non-solvable with modulo trick};
        \node[box, below=of 2] (3) {Develop lemmas for congruence solution sets};
        \node[box, below=of 3] (4) {Theorem 1: $1+(pq)^a = p^b+q^c$ is non-solvable};
        \node[box, below=of 4] (5) {Generalize definitions for other equation forms};
        \node[box, below=of 5] (6) {Lemma 3: Period of exponent residues};
        \node[box, below=of 6] (7) {Theorem 2: $a^x-a^y = b^z-b^w$ is non-solvable};
        
        % Arrows
        \draw[arrow] (1) -- (2);
        \draw[arrow] (2) -- (3);
        \draw[arrow] (3) -- (4);
        \draw[arrow] (4) -- (5);
        \draw[arrow] (5) -- (6);
        \draw[arrow] (6) -- (7);
        
        \end{tikzpicture}}
        \caption{Research methodology flowchart}
    \end{figure}
\end{frame}

% ================================================================================
% Section 6: Results
% ================================================================================
